\documentclass[aspectratio=169]{beamer}
\usetheme{Madrid}
\usecolortheme{default}
\usepackage[utf8]{inputenc}
\usepackage{amsmath}
\usepackage{amsfonts}
\usepackage{amssymb}
\usepackage{enumerate}

\setbeamertemplate{navigation symbols}{}
\setbeamerfont{block body}{size=\small}
\setbeamerfont{alertblock body}{size=\small}

\title{CRISC Domain B -- 100 Real-Life Questions}
\subtitle{IT Risk Assessment -- Practical Scenarios for Daily Challenges}
\author{CRISC Mastery Series}
\date{\today}

% ------------------------------------------------------------
% \hq{number}{topic}{question}{options}{letter}{explanation}{trap}
% Creates one question frame and one explanation frame.
% ------------------------------------------------------------
\newcommand{\hq}[7]{%
	\begin{frame}{Q#1: #2}
		\begin{block}{Question}
			#3
		\end{block}
		\begin{enumerate}[(A)]
			#4
		\end{enumerate}
	\end{frame}
	\begin{frame}{Q#1: Explanation}
		\begin{block}{Correct Answer: #5}
			\begin{itemize}
				#6
			\end{itemize}
		\end{block}
		\begin{alertblock}{Key Takeaway}
			#7
		\end{alertblock}
	\end{frame}
}

\begin{document}
	
	\begin{frame}
		\titlepage
	\end{frame}
	
	\begin{frame}{Instructions}
		\begin{alertblock}{How to Use This Deck}
			Think like a risk practitioner in a real organization.\\
			Focus on business context, ownership, and practical decision-making.\\
			The answer and a key takeaway are on the slide following each question.
		\end{alertblock}
	\end{frame}
	
	\begin{frame}{Outline}
		\tableofcontents
	\end{frame}
	
	% ============================================================
	\section{Risk Identification}
	% ============================================================
	
	\hq{1}{Startup Cloud Migration}
	{A fast-growing startup is moving all its systems to the cloud. The new CISO wants to understand the risks. What should be the \textbf{FIRST} step?}
	{\item Perform a penetration test on the cloud provider
		\item Review the provider's SOC 2 report
		\item Meet with business leaders to understand their goals and priorities
		\item Negotiate a better SLA to cover security incidents}
	{C}
	{\item Risk is about the impact on business objectives.
		\item You must know what the business is trying to achieve to identify what could go wrong.
		\item Technical checks and contract reviews come after understanding the business context.}
	{You can't assess risk without first knowing the business objectives.}
	
	\hq{2}{The Shadow IT Discovery}
	{A security analyst discovers a marketing team is using an unsanctioned SaaS tool to store customer analytics. What is the \textbf{BEST} course of action?}
	{\item Immediately block access to the tool
		\item Report the marketing team to their manager
		\item Ignore it, as the data is only for analytics
		\item Talk to the marketing team to understand how they use the tool and what data is involved}
	{D}
	{\item The first step is to understand the situation.
		\item What data is being stored? What is the business value?
		\item Blocking might disrupt important work and damage security's relationship with the business.
		\item A risk assessment should follow, but discovery and understanding come first.}
	{Understand the "why" before taking disruptive action.}
	
	\hq{3}{New Regulation Impact}
	{A new data privacy regulation is passed that will take effect in 12 months. What should the risk practitioner do \textbf{NOW}?}
	{\item Wait for the effective date to see if it is enforced
		\item Purchase new privacy software
		\item Conduct a gap assessment to see how current practices align with the new law
		\item Update the company website's privacy policy}
	{C}
	{\item The law will create new obligations.
		\item A gap assessment identifies where the organization is currently non-compliant.
		\item This allows for planning, budgeting, and remediation before the deadline.
		\item Actions like buying software or updating policies should follow the assessment.}
	{Understand the gap before you try to close it.}
	
	\hq{4}{Board Risk Report}
	{You are preparing a risk report for the board of directors. Which of the following is the \textbf{MOST} important to include?}
	{\item A list of all 500 vulnerabilities found in the last scan
		\item A technical summary of the firewall rules
		\item A summary of the top risks, their potential business impact, and status versus risk appetite
		\item The number of phishing emails blocked by the email gateway}
	{C}
	{\item The board is responsible for governance and strategic oversight.
		\item They need information that is relevant to business objectives.
		\item Technical details and raw counts are not useful for high-level decision-making.}
	{Speak the board's language: business impact and risk posture.}
	
	\hq{5}{After a Major Incident}
	{A ransomware attack was just contained. The business is back online. What is the \textbf{MOST} important risk-related activity to do next?}
	{\item Buy a new firewall
		\item Have a lessons-learned meeting to update the risk register and scenarios
		\item Fire the employee whose credentials were stolen
		\item Increase the security budget}
	{B}
	{\item The incident was a risk that materialized.
		\item Root cause analysis reveals why controls failed and how the risk assessment missed or underestimated the scenario.
		\item The risk register, scenarios, and controls should be updated to reflect this new information.}
	{Treat a major incident as a learning opportunity to improve the risk process.}
	
	\hq{6}{New E-commerce Platform}
	{The business wants to launch a new e-commerce platform in three months. As the risk practitioner, what is your \textbf{FIRST} priority?}
	{\item Ensure all payment card data is encrypted
		\item Perform a threat modeling exercise for the new platform
		\item Include the risk practitioner in the project from the beginning to identify risks as they arise
		\item Request a budget for a 24/7 security operations center}
	{C}
	{\item Risk identification is a continuous process.
		\item Being involved early allows you to influence the design and identify risks before they become embedded in the solution.
		\item Specific controls and activities should be based on identified risks.}
	{Embed risk management into the project, don't bolt it on at the end.}
	
	\hq{7}{Cloud Provider Assurance}
	{Your critical payroll system is moving to a new SaaS provider. Which of the following sources of assurance is \textbf{MOST} reliable?}
	{\item The provider's marketing materials claiming they are "secure"
		\item A recent, independent, third-party audit report (e.g., SOC 2 Type II)
		\item The provider's answers to a security questionnaire
		\item A verbal assurance from the provider's sales representative}
	{B}
	{\item An independent audit provides objective, evidence-based validation of controls over a period.
		\item Marketing materials and verbal assurances are not verifiable.
		\item Questionnaires are helpful but self-reported by the vendor.}
	{Seek objective, independent evidence for critical vendors.}
	
	\hq{8}{The Delphi Technique}
	{You need to identify the potential risks of a new, innovative technology that has no historical data. You want to avoid one person's strong opinion dominating the discussion. Which technique is \textbf{BEST}?}
	{\item Brainstorming with all stakeholders
		\item The Delphi technique, using anonymous input from experts
		\item A checklist-based review from a past project
		\item A structured walkthrough of the system design}
	{B}
	{\item The Delphi technique uses anonymous, structured questionnaires with experts.
		\item This prevents a single strong personality from unduly influencing the group.
		\item It is excellent for new or complex situations where historical data is lacking.}
	{For unbiased expert consensus, use the Delphi technique.}
	
	\hq{9}{Vendor Risk Assessment}
	{Your company is onboarding a new vendor for a low-risk, non-critical service. What is the \textbf{MOST} appropriate level of risk assessment?}
	{\item No assessment is needed since it's a low-risk service
		\item A full, in-person, week-long audit of the vendor's facilities
		\item A proportionate assessment, such as a security questionnaire, based on the risk
		\item Rely on the vendor's reputation in the market}
	{C}
	{\item Risk assessment should be proportionate to the risk.
		\item A low-risk vendor does not require a resource-intensive audit.
		\item A questionnaire is a reasonable, cost-effective way to gain some assurance.
		\item No assessment or relying only on reputation leaves the organization exposed.}
	{Match the depth of your assessment to the level of risk.}
	
	\hq{10}{Merger and Acquisition}
	{Your company is acquiring another firm. You need to understand the cybersecurity risks of the acquisition. What should you do \textbf{FIRST}?}
	{\item Immediately integrate the acquired company's network into yours
		\item Conduct a risk assessment of the acquired company's environment
		\item Terminate all IT staff at the acquired company
		\item Migrate all their data to your data centers}
	{B}
	{\item Connecting networks or migrating data before an assessment can introduce unknown risks.
		\item An assessment will identify vulnerabilities, control gaps, and areas of non-compliance.
		\item This allows for a planned, risk-based integration.}
	{Assess the risk before you connect the networks.}
	
	\hq{11}{Phishing Simulation Results}
	{A phishing simulation shows a 30\% click rate. What is the \textbf{BEST} way to interpret and act on this result?}
	{\item It's a technical problem; the email filter needs to be improved
		\item It confirms that employees are the weakest link and should be disciplined
		\item It indicates a need for targeted security awareness training
		\item The simulation was too hard; the next one should be easier}
	{C}
	{\item A high click rate is a human risk, not just a technical one.
		\item The most effective response is to improve the knowledge and skills of the employees through training.
		\item While filters are important, they will never be 100\% effective.}
	{Address human risk with education, not just technology.}
	
	\hq{12}{Third-Party Data Breach}
	{You learn that a third-party vendor holding your customer data has suffered a breach. What is your \textbf{FIRST} priority?}
	{\item Terminate the contract with the vendor
		\item Determine the scope of the breach and what data was affected
		\item Issue a press release to get ahead of the story
		\item Sue the vendor for damages}
	{B}
	{\item You cannot respond effectively without knowing what happened.
		\item Understanding the scope and impact is essential for determining your notification obligations (legal, regulatory, customer) and the appropriate response.
		\item Legal and contractual actions follow the initial investigation.}
	{Understand the incident before you react to it.}
	
	\hq{13}{Board Appetite Statement}
	{The board has a very low risk appetite for data breaches. Which of the following activities is \textbf{MOST} aligned with this statement?}
	{\item Investing in the latest marketing automation tool
		\item Prioritizing the implementation of data loss prevention (DLP) controls
		\item Upgrading the office coffee machines
		\item Reducing the security budget to increase profitability}
	{B}
	{\item The risk appetite is the amount of risk the organization is willing to accept.
		\item A low appetite for data breaches means the organization should invest in controls that prevent them.
		\item DLP directly addresses the risk of data loss.
		\item Other activities do not directly relate to the stated appetite.}
	{Investments should be aligned with the board's risk appetite.}
	
	\hq{14}{Risk Register Update}
	{A project manager tells you they "don't have time" to update the risk register. What is the \textbf{BEST} way to respond?}
	{\item Escalate the issue to their manager
		\item Update the register for them
		\item Explain how the register helps them track issues and protects them if a risk materializes
		\item Agree that it is a low priority}
	{C}
	{\item The register is a tool for the project manager, not just a compliance exercise.
		\item It helps them manage their project and provides a record of due diligence.
		\item Framing it as a benefit, rather than a burden, is more likely to gain buy-in.}
	{Help project managers see the value of risk management for themselves.}
	
	\hq{15}{Emerging Technology Adoption}
	{The business wants to adopt a new generative AI tool. What is the \textbf{MOST} important risk consideration?}
	{\item The cost of the tool
		\item The potential for the AI to generate incorrect or biased information
		\item The vendor's stock price
		\item The ease of use for employees}
	{B}
	{\item Generative AI introduces new risks, such as hallucination, bias, and data leakage.
		\item Understanding these potential negative outcomes (incorrect decisions, reputational damage) is a primary risk concern.
		\item While cost is a factor, it is not the primary risk driver.}
	{For emerging tech, focus on the new and unique risks it introduces.}
	
	% ============================================================
	\section{Risk Analysis}
	% ============================================================
	
	\hq{16}{SLE Calculation}
	{An asset is valued at \$200,000. A risk event would destroy 40\% of its value. What is the Single Loss Expectancy (SLE)?}
	{\item \$80,000
		\item \$200,000
		\item \$120,000
		\item \$40,000}
	{A}
	{\item SLE = Asset Value (AV) × Exposure Factor (EF).
		\item SLE = \$200,000 × 0.40 = \$80,000.
		\item This is the expected loss from a single occurrence of the risk event.}
	{SLE is the loss from one event: AV × EF.}
	
	\hq{17}{ALE Calculation}
	{The SLE for a specific risk is \$50,000. The organization expects this event to occur once every two years. What is the Annualized Loss Expectancy (ALE)?}
	{\item \$100,000
		\item \$50,000
		\item \$25,000
		\item \$200,000}
	{C}
	{\item The Annualized Rate of Occurrence (ARO) is 0.5 (once every two years).
		\item ALE = SLE × ARO.
		\item ALE = \$50,000 × 0.5 = \$25,000.}
	{ALE is the expected annual loss: SLE × ARO.}
	
	\hq{18}{Control Cost Justification}
	{A control would cost \$30,000 per year. It reduces the ALE of a risk from \$100,000 to \$20,000. What is the net benefit of the control?}
	{\item \$80,000
		\item \$30,000
		\item \$50,000
		\item -\$10,000}
	{C}
	{\item The control reduces the expected loss by \$80,000 (\$100,000 - \$20,000).
		\item The net benefit is the loss reduction minus the control cost.
		\item Net Benefit = \$80,000 - \$30,000 = \$50,000.}
	{Justify controls by their net benefit, not just their cost.}
	
	\hq{19}{Qualitative vs Quantitative}
	{A manager argues for a quantitative risk analysis for a new project. What is the \textbf{PRIMARY} benefit they are likely seeking?}
	{\item It is faster and cheaper to perform
		\item It provides a monetary value for the risk, aiding in financial decisions
		\item It is based on the subjective opinions of experts
		\item It requires less data to be effective}
	{B}
	{\item Quantitative analysis provides a monetary (e.g., dollar) value for risk.
		\item This is extremely useful for comparing risks, prioritizing investments, and communicating with financial stakeholders (like the CFO or board).
		\item Qualitative analysis is faster and cheaper but provides subjective ratings (High/Medium/Low).}
	{Choose quantitative analysis when you need to speak in dollars.}
	
	\hq{20}{Risk Score Aggregation}
	{You have five independent risks, each with a 'Medium' rating. The CRO asks for an aggregate view. What is the \textbf{MOST} accurate statement?}
	{\item The aggregate risk is 'Medium'
		\item The aggregate risk is five times the individual risk
		\item The aggregate risk is likely higher than any single 'Medium' risk and should be assessed as a portfolio
		\item The aggregate risk is zero because they are independent}
	{C}
	{\item Risk is not additive in a simple way.
		\item Five independent medium risks can create a portfolio of exposure that is higher than any individual risk.
		\item Aggregating risks by business unit, asset, or risk type is a critical part of enterprise risk management.}
	{Look at the portfolio, not just the individual items.}
	
	\hq{21}{Inherent vs Residual Risk}
	{A firewall is installed to protect a web server. The risk of an external attack before the firewall was installed is an example of what type of risk?}
	{\item Residual risk
		\item Current risk
		\item Inherent risk
		\item Transferred risk}
	{C}
	{\item Inherent risk is the level of risk present before any controls are applied.
		\item The firewall is the control.
		\item The risk remaining after the firewall is implemented is the residual risk.}
	{Inherent is before controls; residual is after.}
	
	\hq{22}{Risk Appetite and Tolerance}
	{A risk has been assessed. Its level is currently above the stated risk appetite but below the risk tolerance. What is the \textbf{MOST} appropriate action?}
	{\item Do nothing, as it is below tolerance
		\item Implement immediate corrective controls to bring it within appetite
		\item Monitor the risk closely and consider action if it trends toward tolerance
		\item Ignore the appetite statement and use only tolerance}
	{C}
	{\item Appetite is the desired level of risk.
		\item Tolerance is the acceptable deviation from appetite.
		\item A risk between appetite and tolerance is not in an ideal state but is still within acceptable bounds. It requires monitoring, not immediate panic, but action may be needed if it worsens.}
	{Appetite is the goal, tolerance is the buffer. Manage the space between.}
	
	\hq{23}{New Threat Intelligence}
	{You receive threat intelligence about a new exploit targeting a software your company uses. What is the \textbf{FIRST} step in analyzing this information?}
	{\item Immediately shut down the affected systems
		\item Forward the intelligence to the IT team to patch
		\item Verify if your organization is actually vulnerable to this specific exploit
		\item Wait for a vendor patch to be released}
	{C}
	{\item Threat intelligence is not a confirmed risk until you verify your exposure.
		\item Do you use the specific version of the software? Do you have compensating controls?
		\item This analysis determines if and how you should respond.}
	{Threat intel is a data point; validate its applicability to your environment.}
	
	\hq{24}{Scenario Development}
	{When developing a risk scenario for a new database, what is the \textbf{MOST} important element to include?}
	{\item The version of the database software
		\item The name of the database administrator
		\item The potential threat, the vulnerability it could exploit, and the business impact
		\item The server's IP address}
	{C}
	{\item A risk scenario must link a cause (threat and vulnerability) to a consequence (business impact).
		\item The technical details (version, IP) are inputs, but the scenario is a narrative of what could happen and why it matters.
		\item Without a business impact, it's not a fully formed risk.}
	{A good risk scenario tells a story: cause, event, and consequence.}
	
	\hq{25}{Expert Judgment Bias}
	{A team of senior IT engineers assesses a risk as "low" because the technology is familiar to them. What cognitive bias might be at play?}
	{\item Confirmation bias
		\item Anchoring bias
		\item Availability bias
		\item Overconfidence bias}
	{D}
	{\item Overconfidence bias occurs when experts overestimate their own abilities and the safety of their domain.
		\item Familiarity with the technology can lead to an underestimation of the risks associated with it.
		\item Seeking diverse perspectives can help counter this.}
	{Experts can be overconfident; challenge their assumptions.}
	
	\hq{26}{Key Risk Indicator (KRI)}
	{Which of the following is the \textbf{BEST} example of a leading Key Risk Indicator (KRI) for a system availability risk?}
	{\item The total hours of unplanned downtime last month
		\item The number of production incidents last quarter
		\item The trend of increasing server CPU utilization
		\item The cost of the last outage}
	{C}
	{\item A leading indicator is predictive; it changes before the risk event occurs.
		\item Increasing CPU utilization suggests the system is approaching a threshold that could lead to an outage.
		\item The other options are lagging indicators; they report on events that have already happened.}
	{KRIs should look forward, not backward.}
	
	\hq{27}{Risk Data Quality}
	{You are using historical incident data from three years ago to estimate the likelihood of a risk. The business has since transformed its technology stack. What is the \textbf{PRIMARY} concern with this data?}
	{\item The data is too old.
		\item The data is not relevant to the current environment.
		\item There is not enough data.
		\item The data is not accurate.}
	{B}
	{\item The key issue is the relevance or representativeness of the data.
		\item A transformed technology stack means past data may not reflect current threats, vulnerabilities, or controls.
		\item The analysis could be significantly over- or under-stated as a result.}
	{Old data isn't just old; it may be irrelevant. Check its context.}
	
	\hq{28}{Control Effectiveness}
	{After implementing a new control, you want to determine its effectiveness. What is the \textbf{MOST} direct way to do this?}
	{\item Review the design documents for the control
		\item Interview the control owner
		\item Test the control to see if it is operating as intended
		\item Read the vendor's marketing brochure for the control}
	{C}
	{\item A control's design can look good on paper, but it's only effective if it operates correctly in the real world.
		\item Testing provides direct evidence of the control's operating effectiveness.
		\item The other options are indirect or provide evidence only about the design.}
	{The only way to know if a control works is to test it.}
	
	\hq{29}{Risk Assessment Scope}
	{An organization is performing its annual risk assessment. What is the \textbf{BEST} way to define the scope?}
	{\item All IT systems managed by the IT department
		\item The systems that were tested last year
		\item The critical business processes and the IT systems that support them
		\item All systems mentioned in the latest audit report}
	{C}
	{\item Risk is about impact on the business.
		\item The scope should be driven by the business processes that are most critical to achieving objectives.
		\item This ensures that IT risk management is aligned with enterprise risk management.}
	{Scope your risk assessment around business value, not just IT assets.}
	
	\hq{30}{Residual Risk Acceptance}
	{After implementing controls, the residual risk is still above the organization's risk appetite. Who is the \textbf{BEST} person to formally accept this risk?}
	{\item The control owner
		\item The IT manager
		\item The business owner who benefits from the process
		\item The risk practitioner}
	{C}
	{\item Risk acceptance is a business decision.
		\item The person who owns the business process and bears the consequences of the risk should be the one to accept it.
		\item They have the authority and accountability to do so.}
	{The person who bears the risk should accept the risk.}
	
	% ============================================================
	\section{Risk Evaluation and Response}
	% ============================================================
	
	\hq{31}{Risk Response Options}
	{A risk has been identified. The cost to mitigate it exceeds the potential loss. What is the \textbf{MOST} likely appropriate risk response?}
	{\item Avoid the risk by ceasing the activity
		\item Transfer the risk to an insurance company
		\item Accept the risk and monitor it
		\item Implement the expensive mitigation anyway}
	{C}
	{\item When the cost of a control is greater than the benefit it provides, it is not a sound investment.
		\item Accepting the risk is a valid response when the risk is within tolerance and the cost to treat it is disproportionate.
		\item The acceptance must be documented and approved.}
	{Don't spend more to fix a risk than the risk is worth.}
	
	\hq{32}{Insurance as a Control}
	{A company buys cyber insurance. This is an example of which risk response strategy?}
	{\item Mitigation
		\item Avoidance
		\item Transfer
		\item Acceptance}
	{C}
	{\item Risk transfer involves shifting the financial consequence of a risk to a third party, like an insurer.
		\item It does not eliminate the risk itself, but it transfers the financial impact.
		\item The organization is still responsible for managing the incident and its reputational damage.}
	{Insurance transfers financial impact, not accountability.}
	
	\hq{33}{Risk Mitigation Example}
	{Implementing a multi-factor authentication (MFA) system is an example of which risk response?}
	{\item Risk acceptance
		\item Risk transfer
		\item Risk avoidance
		\item Risk mitigation}
	{D}
	{\item Risk mitigation involves taking action to reduce the likelihood or impact of a risk.
		\item MFA reduces the likelihood of unauthorized access by making credential theft less effective.
		\item It is a direct action to lower the risk.}
	{Mitigation is about reducing the risk, not just shifting it.}
	
	\hq{34}{Risk Avoidance Example}
	{A company decides not to enter a new market because of the high regulatory and legal risks. This is an example of what?}
	{\item Risk acceptance
		\item Risk mitigation
		\item Risk avoidance
		\item Risk transfer}
	{C}
	{\item Risk avoidance is the decision to not perform an activity that carries risk.
		\item By not entering the market, the company completely avoids the associated risks.
		\item This is different from mitigating a risk to make an activity safer.}
	{Avoidance means you stop doing the risky thing entirely.}
	
	\hq{35}{Compensating Controls}
	{A legacy system cannot be patched and has a known vulnerability. To manage the risk, the team places the system behind a firewall and monitors its traffic closely. These new measures are examples of what?}
	{\item Preventive controls
		\item Detective controls
		\item Compensating controls
		\item Corrective controls}
	{C}
	{\item Compensating controls are alternative safeguards implemented when the primary control (patching) is not feasible.
		\item They are designed to reduce the risk to an acceptable level through other means.
		\item A firewall and monitoring act as a substitute for the missing patch.}
	{When you can't fix the root cause, compensate with other controls.}
	
	\hq{36}{Risk Owner Role}
	{What is the primary responsibility of a risk owner?}
	{\item To perform the technical testing of controls
		\item To implement the risk response plan and be accountable for the risk
		\item To write the risk management policy
		\item To conduct audits of the risk management process}
	{B}
	{\item The risk owner is accountable for managing a specific risk.
		\item They are responsible for ensuring that the chosen risk response is implemented and that the risk remains within acceptable limits.
		\item This includes providing resources and making decisions.}
	{The risk owner is the accountable person for a specific risk.}
	
	\hq{37}{Control Owner Role}
	{After a risk response is approved, who is responsible for the day-to-day operation of the implemented control?}
	{\item The risk owner
		\item The control owner
		\item The internal auditor
		\item The CISO}
	{B}
	{\item The risk owner decides *what* to do, but the control owner is responsible for the *implementation and operation* of the specific control.
		\item The control owner is often a technical or operational person who manages the control on an ongoing basis.}
	{The risk owner owns the risk; the control owner owns the control.}
	
	\hq{38}{Responding to an Audit Finding}
	{An internal audit finds that a key control is not operating effectively. What is the \textbf{MOST} appropriate action for the risk practitioner?}
	{\item Ignore the finding as it was an internal audit
		\item Argue that the control is too expensive to fix
		\item Update the risk register and re-assess the risk with the risk owner
		\item Wait for the external audit to see if they find the same issue}
	{C}
	{\item An audit finding provides new evidence about a control's weakness.
		\item This means the risk it was designed to mitigate is now higher than previously thought.
		\item The risk register, risk assessment, and response plan must be updated to reflect this new reality.}
	{An audit finding is a new piece of risk information. Act on it.}
	
	\hq{39}{Risk Monitoring}
	{What is the primary purpose of continuous risk monitoring?}
	{\item To satisfy regulatory requirements
		\item To detect changes in the risk environment and risk profile
		\item To justify the budget for security tools
		\item To generate data for reports}
	{B}
	{\item The risk environment is constantly changing (new threats, business changes, control drift).
		\item Monitoring provides timely awareness of these changes so that the organization can adapt its risk management activities proactively.
		\item Compliance and reporting are secondary benefits.}
	{Monitor risk to know when the picture changes, not just for reports.}
	
	\hq{40}{Risk Reporting to Business}
	{When communicating a risk to a business manager, what is the \textbf{MOST} effective way to frame it?}
	{\item Describe the vulnerability in technical terms, including the CVE number
		\item Provide a list of the specific firewall rules that were bypassed
		\item Explain the potential impact on their business goals, such as revenue, customer trust, or operational efficiency
		\item Give them a copy of the technical risk assessment report}
	{C}
	{\item Business managers care about business objectives.
		\item Framing the risk in terms of its potential impact on their area (e.g., "this could disrupt our sales process," "this could lead to fines") is the most effective way to gain their attention and support for a response.}
	{Translate technical risk into business impact.}
	
	\hq{41}{New Business Initiative}
	{The business wants to launch a new service that uses a new technology. The risk assessment shows the inherent risk is high. What should the risk practitioner do?}
	{\item Block the initiative because the risk is too high
		\item Approve the initiative without further discussion
		\item Work with the business to identify controls to mitigate the risk to an acceptable level
		\item Tell the business they must accept all the risk}
	{C}
	{\item The role of risk management is to enable the business, not to be a roadblock.
		\item High inherent risk does not mean an activity should be stopped; it means controls are needed to reduce the risk.
		\item The practitioner should collaborate with the business to find a way to proceed safely.}
	{Your job is to help the business take smart risks, not avoid all risks.}
	
	\hq{42}{Key Performance Indicator (KPI)}
	{Which of the following is a Key Performance Indicator (KPI) for a vulnerability management program?}
	{\item The number of high-severity vulnerabilities currently open
		\item The percentage of critical vulnerabilities remediated within the defined 30-day SLA
		\item The number of new exploits in the wild this quarter
		\item The total number of systems in the environment}
	{B}
	{\item KPIs measure the performance and effectiveness of a process or control.
		\item "Percentage remediated within SLA" directly measures how well the vulnerability management process is performing against its goal.
		\item The other options are better classified as risk indicators (KRIs) or simple metrics.}
	{KPIs measure performance against a goal or SLA.}
	
	\hq{43}{Key Control Indicator (KCI)}
	{Which of the following is a Key Control Indicator (KCI)?}
	{\item The number of phishing emails that bypassed the email gateway
		\item The percentage of servers that have the latest security patches installed
		\item The trend of external attacks on the organization's network
		\item The number of employees who failed the last phishing simulation}
	{B}
	{\item A KCI measures whether a specific control is in place and operating as designed.
		\item The "percentage of servers with the latest security patches" directly measures the performance of the patch management control.
		\item The other options are more related to risk (KRIs) or performance (KPIs).}
	{KCIs measure the control itself.}
	
	\hq{44}{Reporting an Accepted Risk}
	{A business unit has formally accepted a risk. What should the risk practitioner do with this information?}
	{\item Forget about it, the risk is now their problem
		\item Escalate it to the board
		\item Document it in the risk register and continue to monitor for any changes
		\item Try to secretly implement controls anyway}
	{C}
	{\item Accepting a risk is a formal decision.
		\item The decision and its rationale must be documented (in the risk register) for accountability and future reference.
		\item The risk practitioner should continue to monitor the risk to ensure the basis for acceptance remains valid.}
	{Accepted risk is still tracked risk. Document and monitor it.}
	
	\hq{45}{Risk Assessment Frequency}
	{How often should a risk assessment be performed?}
	{\item Only once per year
		\item Only when a major incident occurs
		\item At planned intervals and whenever there is a significant change
		\item Every time a new employee is hired}
	{C}
	{\item Risk management is continuous, but formal risk assessments are periodic.
		\item They should be performed on a regular schedule (e.g., annually) to catch drift.
		\item They must also be triggered by significant events, such as a merger, a new technology adoption, or a major incident.}
	{Periodic and event-driven: the two triggers for a risk assessment.}
	
	\hq{46}{Risk Appetite Communication}
	{How can an organization ensure that its risk appetite is understood and used by all employees?}
	{\item By mentioning it once in the annual code of conduct training
		\item By translating the high-level appetite statement into specific policies, procedures, and decision-making criteria for different teams
		\item By posting the statement on the company intranet
		\item By having the CRO sign off on the statement}
	{B}
	{\item A vague, high-level statement is not actionable for most employees.
		\item To be effective, the appetite must be operationalized—translated into concrete rules like "any risk above \$X requires CFO approval" or "we will not use any software that is not on the approved list."
		\item This makes it part of daily work.}
	{Make risk appetite practical, not just a philosophical statement.}
	
	\hq{47}{Mitigation vs. Acceptance Trade-off}
	{A risk has a potential loss of \$1 million. A control to mitigate it would cost \$200,000 per year. The risk is considered to have a low likelihood. What is the \textbf{MOST} important factor to consider before deciding?}
	{\item The \$200,000 cost is only 20\% of the loss.
		\item The precise likelihood of the risk occurring.
		\item The risk practitioner's opinion.
		\item What competitors are doing about this risk.}
	{B}
	{\item You have the impact (\$1M) and the cost (\$200k).
		\item The likelihood is the missing piece of the equation.
		\item If the risk is extremely rare (e.g., 1-in-100 years), the annualized loss expectancy might be very low, making the \$200k/year control difficult to justify. The final decision hinges on a better estimate of the likelihood.}
	{To justify a control, you need to know the likelihood, not just the impact.}
	
	\hq{48}{Shadow AI}
	{A software developer is using a public generative AI tool to write code, copy-pasting company code into it. What is the \textbf{PRIMARY} risk?}
	{\item The AI might generate buggy code.
		\item The developer might become too dependent on the AI.
		\item Company intellectual property could be leaked into the public domain.
		\item The AI tool might be unavailable.}
	{C}
	{\item Pasting proprietary source code into public tools can result in that code being retained, used for training, or otherwise exposed.
		\item This is a significant data leak and intellectual property risk.
		\item While code quality is a concern, the confidentiality risk is often the most severe.}
	{Public AI tools are not a safe place for proprietary data.}
	
	\hq{49}{Vendor Lock-in}
	{A company has all its data in a proprietary cloud platform, making it very difficult to move to another provider. What type of risk does this represent?}
	{\item Compliance risk
		\item Concentration risk
		\item Reputational risk
		\item Strategic risk}
	{B}
	{\item Concentration risk is the risk of relying on a single partner or technology.
		\item Being locked in increases the impact of any failure, price increase, or policy change by that single vendor.
		\item It limits the organization's flexibility and negotiating power.}
	{Being too dependent on one provider is a risk in itself.}
	
	\hq{50}{Post-Incident KRI Update}
	{After a major outage caused by a database corruption, what is the \textbf{BEST} way to update KRIs?}
	{\item No update is needed, the incident is over.
		\item Add a KRI for the "number of database backups completed."
		\item Add a KRI for "database integrity check failure rate" or "backup restore test failure rate."
		\item Remove all existing KRIs related to databases.}
	{C}
	{\item The KRIs should have provided an early warning of this incident.
		\item A database integrity KRI or a backup restore test KRI might have indicated a problem before the outage.
		\item The goal is to add leading indicators that would have helped predict or prevent the specific event that occurred.}
	{After an incident, ask: "What indicator would have warned us?" and add it.}
	
	% ============================================================
	\section{Business Impact Analysis (BIA)}
	% ============================================================
	
	\hq{51}{BIA Purpose}
	{What is the primary purpose of a Business Impact Analysis (BIA)?}
	{\item To list all the vulnerabilities in the IT environment
		\item To determine the Recovery Time Objective (RTO) and Recovery Point Objective (RPO) for critical business processes
		\item To calculate the financial return on investment for security tools
		\item To document the IT architecture}
	{B}
	{\item The BIA identifies which business processes are most critical and what the impact of a disruption would be.
		\item The key outputs of a BIA are the recovery objectives: RTO (how quickly must we recover?) and RPO (how much data can we lose?).
		\item It is a business-driven process.}
	{BIA's main job is to set RTO and RPO for what matters most.}
	
	\hq{52}{RTO vs RPO}
	{A system can tolerate 4 hours of downtime and a maximum of 15 minutes of data loss. What are its RTO and RPO, respectively?}
	{\item RTO = 15 minutes, RPO = 4 hours
		\item RTO = 4 hours, RPO = 15 minutes
		\item RTO = 4 hours, RPO = 4 hours
		\item RTO = 15 minutes, RPO = 15 minutes}
	{B}
	{\item Recovery Time Objective (RTO) is the maximum acceptable downtime. So, RTO = 4 hours.
		\item Recovery Point Objective (RPO) is the maximum acceptable data loss, measured in time. So, RPO = 15 minutes.
		\item Think RTO = Time to recover; RPO = Point in time of last good data.}
	{RTO is about time down; RPO is about data lost.}
	
	\hq{53}{BIA Dependency}
	{During a BIA, it is discovered that the payroll system (RTO 24h) depends on the HR database (RTO 48h). What is the \textbf{PRIMARY} issue?}
	{\item There is no issue, they are different systems.
		\item The HR database has a higher RTO than the system it supports, creating a dependency conflict.
		\item The payroll system's RTO is too strict.
		\item The HR database should not have an RTO.}
	{B}
	{\item A system cannot be recovered before the systems it depends on.
		\item If the payroll system must be back in 24 hours, but its dependency (the HR database) can be down for 48 hours, the payroll system's RTO is impossible to meet.
		\item Dependency RTOs must be less than or equal to the RTOs of the systems that depend on them.}
	{A system's RTO is constrained by the RTOs of its dependencies.}
	
	\hq{54}{MTD Definition}
	{Which of the following best describes the Maximum Tolerable Downtime (MTD)?}
	{\item The average time it takes to recover a system
		\item The total amount of time a business process can be unavailable before the impact becomes unacceptable
		\item The time it takes to restore data from backup
		\item The target recovery time for the IT team}
	{B}
	{\item MTD is a business-level concept.
		\item It is the absolute limit—the point beyond which the business faces severe, potentially catastrophic consequences (e.g., bankruptcy, massive regulatory fines).
		\item The RTO should always be less than the MTD.}
	{MTD is the point of no return for the business.}
	
	\hq{55}{BIA Interdependency}
	{A BIA reveals that an e-commerce website depends on three systems: a web server, a database, and a payment gateway. All have an RTO of 2 hours. How should recovery be prioritized?}
	{\item Alphabetically by system name
		\item By the cost of each system
		\item The foundational systems (database, payment gateway) must be recovered before or with the web server for it to function
		\item They can all be recovered simultaneously without issue}
	{C}
	{\item The web server is useless without the database and payment gateway.
		\item The recovery plan must respect these dependencies.
		\item The foundational systems must be brought online first or at the same time to meet the 2-hour RTO for the overall service.}
	{Recover the foundation first, then build on top of it.}
	
	\hq{56}{RPO and Backups}
	{A system has an RPO of 1 hour. The current backup schedule is a full backup nightly at midnight. Is this adequate?}
	{\item Yes, a nightly backup is a standard practice.
		\item No, the backup frequency is too low. The organization could lose up to 24 hours of data, which far exceeds the 1-hour RPO.
		\item Yes, as long as the backups are encrypted.
		\item No, the backups should be weekly instead of nightly.}
	{B}
	{\item The RPO sets the maximum acceptable data loss.
		\item A nightly backup means that if a failure occurs at 11:00 PM, the last backup was at midnight, so nearly 23 hours of data would be lost.
		\item This violates the 1-hour RPO. More frequent backups (or log shipping/continuous replication) are needed.}
	{Your backup schedule must be aligned with your RPO.}
	
	\hq{57}{DR Test Failure}
	{A disaster recovery test reveals that a critical system took 10 hours to restore, but its RTO is 4 hours. What should the risk practitioner do \textbf{FIRST}?}
	{\item Change the RTO to 10 hours
		\item Blame the recovery team
		\item Document the finding and initiate a root cause analysis to improve the recovery process
		\item Do nothing, it was just a test}
	{C}
	{\item A failed DR test is a critical finding. It shows the recovery capability is not as expected.
		\item The first step is to document the gap and find out *why* it took so long.
		\item The RTO should only be changed if the business agrees that a longer downtime is acceptable, which is a separate decision.}
	{A failed test is a real finding. Analyze the cause and fix the plan.}
	
	\hq{58}{BIA and Risk Assessment}
	{How does a Business Impact Analysis (BIA) relate to a risk assessment?}
	{\item They are the same thing
		\item A BIA focuses on the "what" (what processes are critical), while a risk assessment focuses on the "why" (what could disrupt them)
		\item A risk assessment is a subset of a BIA
		\item A BIA is only for natural disasters, while risk assessment is for cyber threats}
	{B}
	{\item They are complementary.
		\item The BIA identifies what is critical and how quickly it needs to be recovered (the "what").
		\item The risk assessment identifies the threats and vulnerabilities that could cause a disruption (the "why" or "how").
		\item You need both to build a resilient business.}
	{BIA tells you what to protect; risk assessment tells you what to protect it from.}
	
	\hq{59}{Recovery Strategy Selection}
	{A critical system has an RTO of 1 hour and an MTD of 4 hours. What is the \textbf{MOST} appropriate recovery strategy?}
	{\item A cold site with data backups
		\item A warm site with data replication
		\item A hot site with real-time synchronous replication
		\item Daily backups to tape stored offsite}
	{C}
	{\item An RTO of 1 hour is extremely aggressive.
		\item Cold, warm, or tape-based strategies will likely take many hours or days to restore.
		\item A hot site with synchronous replication can provide near-instantaneous failover, which is necessary to meet such a tight RTO.}
	{The tighter the RTO, the more you need to invest in a hot, active-active or active-passive recovery solution.}
	
	\hq{60}{Data Classification in BIA}
	{When performing a BIA, why is data classification important?}
	{\item It helps determine the cost of storage.
		\item It helps prioritize which data needs the most stringent recovery objectives (RPO).
		\item It is required for compliance.
		\item It helps the IT team decide which data to delete.}
	{B}
	{\item Highly sensitive or critical data (e.g., customer PII, financial records) may have a much lower RPO (less data loss allowed) than less critical data.
		\item Classification helps ensure that the recovery strategy for different data sets is aligned with its business value and regulatory requirements.}
	{Data classification is a key input to setting RPOs for different systems.}
	
	% ============================================================
	\section{Governance and Risk Culture}
	% ============================================================
	
	\hq{61}{Board Oversight}
	{What is the board's role in IT risk management?}
	{\item To approve all firewall rule changes
		\item To provide oversight, set risk appetite, and ensure that an effective risk management framework exists
		\item To perform the annual penetration test
		\item To manage the day-to-day operations of the security team}
	{B}
	{\item The board is responsible for governance, not management.
		\item They set the tone at the top by defining the risk appetite and ensuring that management has the resources and framework to manage risk effectively.
		\item They provide oversight but do not perform operational tasks.}
	{The board governs; management manages.}
	
	\hq{62}{Three Lines of Defense}
	{In the Three Lines of Defense model, who is responsible for owning and managing risk on a day-to-day basis?}
	{\item The first line (operational management)
		\item The second line (risk and compliance functions)
		\item The third line (internal audit)
		\item The external auditor}
	{A}
	{\item The first line of defense is the operational management (the business itself).
		\item They own the risk and are responsible for implementing and operating the controls.
		\item The second line oversees and advises; the third line provides independent assurance.}
	{The business is the first line of defense; they own the risk.}
	
	\hq{63}{Risk Culture}
	{Which of the following is the \textbf{MOST} effective way to build a strong risk culture?}
	{\item Publishing a risk policy on the company intranet
		\item Giving employees a bonus for never reporting any risks
		\item Leaders consistently modeling the desired risk-aware behavior and rewarding transparency
		\item Requiring all employees to take an annual 30-minute training session}
	{C}
	{\item Culture is shaped by what leaders do, not just what they say.
		\item When leaders consistently demonstrate that risk management is a priority and that they will not punish people for raising concerns, it sends a powerful message.
		\item Training and policies are important, but they are less impactful than consistent leadership behavior.}
	{Culture is caught, not just taught. Leaders set the example.}
	
	\hq{64}{Risk Appetite Statement}
	{Which of the following is a well-defined risk appetite statement?}
	{\item "We will not accept any risk."
		\item "We have a moderate appetite for IT risk."
		\item "We will not tolerate any unauthorized access to customer PII. The loss of more than 1,000 customer records is a significant event requiring board notification."
		\item "We will comply with all laws and regulations."
		{C}
		{\item A good risk appetite statement is specific and measurable.
			\item It defines concrete thresholds that tell you when a risk is approaching a level that requires action or escalation.
			\item Vague terms like "moderate" or "no risk" are not actionable.}
		{A good appetite statement has numbers and consequences.}
		
		\hq{65}{Communication with Regulators}
		{A data breach affecting customer data has occurred. A regulator requests information. What is the \textbf{BEST} approach?}
		{\item Provide the minimum information possible to delay the process
			\item Be transparent and cooperative, providing accurate and timely information as required by law
			\item Give them a copy of the marketing report instead
			\item Only communicate through your lawyers}
		{B}
		{\item Transparency and cooperation are crucial during a regulatory investigation.
			\item Delays or misinformation can lead to penalties and increased scrutiny.
			\item The organization has legal and regulatory obligations to communicate truthfully.}
		{With regulators, honesty and speed are the best policy.}
		
		\hq{66}{Policies vs Standards}
		{What is the difference between a policy and a standard?}
		{\item They are the same thing.
			\item A policy is a high-level statement of intent, while a standard is a mandatory, specific requirement to meet that intent.
			\item A standard is optional, while a policy is mandatory.
			\item A policy is technical, while a standard is for business people.}
		{B}
		{\item Policies are high-level: "All customer data must be encrypted."
			\item Standards are specific and mandatory: "All customer data at rest must be encrypted using AES-256."
			\item Policies are the "what," standards are the "how."}
		{Policies are the "why," standards are the "how much."}
		
		\hq{67}{Exception Management}
		{A business unit requests an exception to a security policy to meet a critical deadline. What is the \textbf{BEST} practice?}
		{\item Grant the exception verbally to save time
			\item Deny the exception; policies are absolute
			\item Follow a formal process to evaluate the risk, approve the exception, document it with an expiration date, and implement compensating controls
			\item Grant the exception permanently so they don't have to ask again}
		{C}
		{\item Exceptions are sometimes necessary, but they must be managed.
			\item A formal process ensures that the risk is understood and approved by the right person.
			\item The exception should be time-bound and include compensating controls to manage the added risk.}
		{Exceptions are okay, but they must be documented, temporary, and controlled.}
		
		\hq{68}{Risk Assessment Scope}
		{An organization is performing its annual risk assessment. What is the \textbf{BEST} way to define the scope?}
		{\item All IT systems managed by the IT department
			\item The systems that were tested last year
			\item The critical business processes and the IT systems that support them
			\item All systems mentioned in the latest audit report}
		{C}
		{\item Risk is about impact on the business.
			\item The scope should be driven by the business processes that are most critical to achieving objectives.
			\item This ensures that IT risk management is aligned with enterprise risk management.}
		{Scope your risk assessment around business value, not just IT assets.}
		
		\hq{69}{Residual Risk Acceptance}
		{After implementing controls, the residual risk is still above the organization's risk appetite. Who is the \textbf{BEST} person to formally accept this risk?}
		{\item The control owner
			\item The IT manager
			\item The business owner who benefits from the process
			\item The risk practitioner}
		{C}
		{\item Risk acceptance is a business decision.
			\item The person who owns the business process and bears the consequences of the risk should be the one to accept it.
			\item They have the authority and accountability to do so.}
		{The person who bears the risk should accept the risk.}
		
		\hq{70}{Third-Party Contract Clauses}
		{To manage third-party risk effectively, which of the following contract clauses is \textbf{MOST} important?}
		{\item A clause allowing the vendor to change their security practices at any time
			\item A clause requiring the vendor to pay for any data breach, regardless of cause
			\item A right-to-audit clause and a requirement to report security incidents within a specific timeframe (e.g., 24 hours)
			\ditem A clause limiting the vendor's liability to the cost of the contract}
		{C}
		{\item A right-to-audit clause allows you to verify the vendor's security claims.
			\item A notification requirement ensures you are aware of an incident quickly so you can respond.
			\item These contractual rights are essential for ongoing oversight.
			\item Unlimited liability clauses are rarely enforceable, and unlimited freedom to change security practices is dangerous.}
		{Oversight doesn't end at onboarding; the contract must enable monitoring.}
		
		\hq{71}{Risk Culture Metric}
		{Which of the following is the \textbf{BEST} indicator of a positive risk culture?}
		{\item The number of risk assessments performed per year
			\item The number of security policies in place
			\item The trend of employees reporting near-misses or potential risks without fear of blame
			\item The amount of money spent on security tools}
		{C}
		{\item A positive risk culture is one where people are empowered and encouraged to speak up about potential problems.
			\item An increasing trend in near-miss reporting is a strong sign of trust and psychological safety.
			\item The other options measure activity or spending, not culture.}
		{Culture is about speaking up, not just ticking boxes.}
		
		\hq{72}{Risk Reporting to Board}
		{What should a risk report to the board of directors focus on?}
		{\item Detailed technical logs of security events
			\item A list of every risk identified in the organization
			\item An overview of the most significant risks, their potential business impact, and the effectiveness of the response
			\item The number of vulnerabilities remediated last month}
		{C}
		{\item The board needs a strategic view, not an operational data dump.
			\item They need to understand the key risks to the business and whether the organization is managing them effectively.
			\item Reports should be concise, focused on impact, and linked to business objectives.}
		{Board reports should be strategic, not technical.}
		
		\hq{73}{Risk Assessment Methodology}
		{An organization is choosing a risk assessment methodology. What is the \textbf{MOST} important consideration?}
		{\item The methodology is the most popular one in the industry
			\item The methodology is easy to use
			\item The methodology aligns with the organization's culture, maturity, and objectives
			\item The methodology requires the least amount of effort}
		{C}
		{\item The "best" methodology is the one that works for the organization.
			\item A highly complex, quantitative methodology may fail in an organization with a less mature risk culture.
			\item The choice should be driven by what is practical, useful, and understood by stakeholders.}
		{Choose a risk methodology that fits your organization, not one that is fashionable.}
		
		\hq{74}{Risk Response in a Project}
		{A project manager is developing a plan to respond to a risk that the project will be delayed by bad weather. They decide to add a buffer to the project schedule. This is an example of what?}
		{\item Risk avoidance
			\item Risk transfer
			\item Risk mitigation
			\ditem Risk acceptance}
		{C}
		{\item The project manager is taking action to reduce the impact of the risk (schedule delay) by adding a buffer.
			\item This is a form of mitigation.
			\item Avoidance would be to not do the project, and transfer would be to insure against the delay.}
		{Mitigation is about taking action to lessen the impact or likelihood.}
		
		\hq{75}{Business Continuity vs Disaster Recovery}
		{What is the difference between Business Continuity (BC) and Disaster Recovery (DR)?}
		{\item They are the same thing.
			\item DR focuses on restoring IT systems after a disaster, while BC focuses on keeping critical business functions running during and after a disaster.
			\item BC focuses on IT, while DR focuses on business processes.
			\item DR is proactive, while BC is reactive.}
		{B}
		{\item DR is a subset of BC, focused specifically on technology.
			\item BC is a broader, business-level plan that covers people, processes, facilities, and technology.
			\item The goal of BC is to keep the business alive, while DR is about getting the IT systems back online.}
		{DR gets IT back up; BC keeps the whole business running.}
		
		% ============================================================
		\section{Monitoring and Reporting}
		% ============================================================
		
		\hq{76}{KRI Design}
		{What makes a good Key Risk Indicator (KRI)?}
		{\item It is easy to measure, even if it's not relevant to the risk
			\item It is predictive of a future risk event and is measurable
			\item It reports on past events in great detail
			\item It is based on a single data point}
		{B}
		{\item A good KRI should be forward-looking (predictive) so it provides an early warning.
			\item It must also be measurable and relevant to a specific risk.
			\item Lagging indicators (reporting on past events) are less useful for proactive management.}
		{KRIs are forward-looking, not backward-looking.}
		
		\hq{77}{KRI Alert}
		{A KRI has breached its threshold. What is the \textbf{FIRST} thing the risk practitioner should do?}
		{\item Immediately implement a new control
			\item Ignore it and wait for the next cycle
			\item Investigate the cause of the breach to understand what is happening
			\ditem Immediately escalate to the board}
		{C}
		{\item A KRI breach is a signal, not a final diagnosis.
			\item The first step is to investigate: Is the data accurate? What is the root cause?
			\item Once the cause is understood, you can determine the appropriate action (which may involve escalating or implementing a control).}
		{When a KRI triggers, investigate before you act.}
		
		\hq{78}{Reporting Metrics}
		{What is the primary difference between a KPI and a KRI?}
		{\item A KPI is for the board, while a KRI is for the IT team.
			\item A KPI measures performance against a goal (e.g., how well a process is working), while a KRI measures exposure to risk.
			\item A KRI is always quantitative, while a KPI is qualitative.
			\item There is no difference; they are interchangeable.}
		{B}
		{\item KPIs (Key Performance Indicators) measure how well you are doing something (e.g., 99.99% uptime).
			\item KRIs (Key Risk Indicators) measure how much risk you are exposed to (e.g., the number of unpatched critical vulnerabilities).
			\item They answer different questions: "How are we doing?" vs. "How much danger are we in?"}
		{KPIs measure performance; KRIs measure risk exposure.}
		
		\hq{79}{Risk Register Contents}
		{Which of the following is the \textbf{MOST} essential piece of information in a risk register entry?}
		{\item The technical name of the vulnerability
			\item The name of the person who identified the risk
			\item The risk statement, the assigned risk owner, and the current risk rating
			\ditem The date the risk will be fixed}
		{C}
		{\item The risk register is a tool for managing risk.
			\item To be actionable, it must state the risk clearly, have an accountable owner, and a current rating to prioritize action.
			\item The other details are helpful but secondary.}
		{A risk without an owner is a problem, not a managed risk.}
		
		\hq{80}{Risk Review with Management}
		{A risk practitioner is preparing to discuss an elevated risk with a business manager. What is the \textbf{MOST} effective approach?}
		{\item Present the technical details of the vulnerability
			\item Show the manager a report with 50 other risks
			\item Tell the manager what control they need to implement
			\item Explain the potential business impact of the risk and discuss options for response}
		{D}
		{\item The goal is to enable a business decision.
			\item Explaining the business impact (e.g., "this could disrupt our order processing") makes the risk relevant.
			\item Presenting options empowers the risk owner to make an informed decision.}
		{Focus on business impact and options, not just technical problems.}
		
		\hq{81}{Audit Findings}
		{An external audit has identified that a control is not operating effectively. What is the risk practitioner's \textbf{MOST} appropriate response?}
		{\item Dispute the finding to protect the company's image
			\item Ignore the finding because it is external
			\item Use the finding as input to update the risk register and re-assess the risk
			\ditem Implement a new control without further analysis}
		{C}
		{\item An audit finding is a valuable piece of evidence about a control weakness.
			\item It means the residual risk is higher than previously thought.
			\item The risk assessment and register must be updated, and a new response plan should be developed with the risk owner.}
		{An audit finding is a gift; it's a new piece of risk information. Use it.}
		
		\hq{82}{Emerging Risks}
		{What is the \textbf{BEST} way to identify emerging risks?}
		{\item Reviewing historical incident data from the last 5 years
			\item Continuously scanning the external environment (threat intelligence, regulatory news, industry trends)
			\item Performing an annual risk assessment
			\ditem Reviewing the most recent audit report}
		{B}
		{\item Emerging risks, by definition, have no history.
			\item They must be identified by looking outward and forward at the changing environment.
			\item This requires continuous monitoring, not just periodic reviews of internal data.}
		{Look outward and forward to see what's coming, not just backward.}
		
		\hq{83}{Key Control Indicator (KCI) Alert}
		{A KCI (e.g., "percentage of access reviews completed") falls to 50\%. What does this suggest?}
		{\item A malicious insider attack is underway
			\item The control designed to review access rights is not being performed consistently
			\item The risk of unauthorized access has decreased
			\ditem Nothing, the KCI is not a useful metric}
		{B}
		{\item A KCI measures the performance of a control.
			\item A drop to 50\% means the access review control is not being done as required.
			\item This means the risk of unauthorized access (the risk the control was meant to mitigate) has likely increased.}
		{A KCI measures the control; a failing KCI means a failing control.}
		
		\hq{84}{Risk Reporting to a Project Team}
		{A project manager needs to report on risk to the project sponsor. What is the \textbf{MOST} important information to include?}
		{\item The number of vulnerabilities found in the latest scan
			\item The top risks that could prevent the project from meeting its objectives, and the planned responses
			\item A detailed list of every risk in the risk register
			\ditem The names of the team members who identified the risks}
		{B}
		{\item The sponsor cares about the project's success.
			\item They need to know what the major threats to success are and what is being done about them.
			\item A focus on the risks that affect project objectives is the most relevant and actionable information.}
		{Project reports should focus on risks to project objectives.}
		
		\hq{85}{Risk and Change Management}
		{A developer wants to make a change to a production system. Why is a formal change management process important from a risk perspective?}
		{\item To slow down the developer
			\item To ensure that the change is made without any risk
			\item To assess the risk of the change, get approval, and have a rollback plan before implementation
			\ditem To satisfy the IT department's need for paperwork}
		{C}
		{\item Changes to production systems can introduce new risks (e.g., outages, new vulnerabilities).
			\item The change management process is a control to manage these risks.
			\item It ensures that changes are tested, approved, and can be reversed if something goes wrong.}
		{Change management is a key control to manage the risk of introducing errors into production.}
		
		% ============================================================
		\section{Conclusion}
		% ============================================================
		
		\begin{frame}{Exam Mindset Summary}
			\begin{block}{Remember These Core Principles}
				\begin{itemize}
					\item \textbf{Start with Business Objectives:} Always understand the business goals before identifying risks.
					\item \textbf{Risk is Dynamic:} Risk is not a one-time assessment; it is a continuous process of identification, assessment, response, and monitoring.
					\item \textbf{Owners Decide:} The risk owner, who is accountable for the business process, is the one who accepts the risk.
					\item \textbf{Assess Before Acting:} Understand the situation and context before taking action, especially with incidents or new technologies.
					\item \textbf{Proportion is Key:} The level of risk assessment and the cost of controls should be commensurate with the level of risk.
					\item \textbf{Communicate in Business Terms:} Translate technical risk into business impact to get stakeholder buy-in.
					\item \textbf{Accountability Cannot be Outsourced:} You can transfer financial risk, but you cannot transfer the accountability for a business process.
				\end{itemize}
			\end{block}
		\end{frame}
		
		\begin{frame}{Thank You}
			\begin{center}
				{\Huge \textbf{Good Luck on Your CRISC Exam!}}
				
				\vspace{1cm}
				
				\normalsize
				The questions in the real exam will test your ability to apply these concepts in practical, real-world scenarios.
			\end{center}
		\end{frame}
		
	\end{document}